\section{Introduction}

Noise-induced transitions are a natural object of study in stochastic dynamics:
random fluctuations can move a trajectory between regions of state space even
when the underlying dynamical rule is fixed. From a statistical point of view,
the question is how accurately such transition times can be detected from an
observed series. Standard change point detection (CPD) simulations usually test
externally inserted shifts in mean, variance, or distribution; they do not cover
events generated by the stochastic process itself.

This paper studies that complementary setting. We generate trajectories from a
nonlinear stochastic dynamical system and label events by a fixed crossing rule
between state-space regions. The generating mechanism is not changed at the
event time; the label marks a transition of the trajectory between regions.
Thus the benchmark is designed for transition-time detection in stochastic
dynamics, not for structural-break detection in the usual parameter-change
sense. Rather than introducing a new detector or escape-rate theory, we use this
dynamics-generated process to compare existing detectors under fixed training,
calibration, testing, and matching rules.

The study makes four contributions. It defines a reproducible transition-time
benchmark with a specified generator, level-crossing labels, metadata, and
independent train/calibration/test splits. It evaluates classical tests,
segmentation and window-based procedures, and learned detectors under the same
matching protocol. It reports \(D\)-dependent rankings, including the
tolerance-sensitive \(D=1.5\) case, and uses a same-noise Transformer diagnostic
to delimit colored-noise neural-advantage claims.

